%lezione 1, 04/03/2026
\chapter{Postulati della Meccanica Quantistica}
\section{Introduzione e costruzione della meccanica classica}
In questa prima parte del corso si formalizzano alcuni concetti già visti
nel corso di Meccanica Quantistica I, partendo dai postulati fondamentali su cui si basa
la teoria.

Per costruire la teoria della meccanica quantistica è necessario basarsi sulle risposte
alle seguenti domande:

\noindent\fbox{%
    \parbox{\linewidth}{
\begin{enumerate}
	    \item \textbf{Come si può descrivere lo stato di un sistema fisico a un tempo
		$ \mathbf{t=t_0} $ fissato?}
	    \item\textbf{Fissato lo stato, come si può predire il risultato della misura di una
		sua osservabile fisica?}
	    \item \textbf{Noto lo stato a $\mathbf{t=t_0}$, come si può determinare lo stato del sistema
		a $\mathbf{t > t_0} $?}
	\end{enumerate}
    }%
}%

Per capire perché queste tre domande bastano a definire la teoria, si considerano prima
le risposte della meccanica classica (nel formalismo hamiltoniano):
\begin{enumerate}
    \item Per un sistema a $ n $ gradi di libertà (dof), lo stato del sistema a $ t=t_0 $
	è definito da $ 2n $ coordinate nello spazio delle fasi
	$ \{ p_0^{\alpha}, q_0^{\alpha} \} $;
    \item Ogni osservabile fisica $ Q $ è funzione reale delle coordinate canoniche
	$ Q = F(q^{\alpha}, p^{\alpha}) $, il valore della funzione nel punto dello
	spazio delle fasi è il valore della osservabile;
    \item L'evoluzione temporale è governata dalle equazioni di Hamilton, che hanno come
	soluzioni curve nello spazio delle fasi dipendenti dallo stato in $ t = t_0 $:
	$ q^{\alpha}(t,q_0^{\alpha},p_0^{\alpha})$, $p^{\alpha}(t,q_0^{\alpha},p_0^{\alpha})$ 
\end{enumerate}
Sostituendo lo spazio delle fasi con quello delle velocità, si definisce analogamente la
teoria nel formalismo lagrangiano.

\section{Postulato 0 della Meccanica Quantistica}

\noindent\fbox{
    \parbox{\textwidth}{
	\textbf{Lo stato di un sistema fisico a $\mathbf{t=t_0}$ è rappresentato da un raggio
	    nello spazio di Hilbert separabile associato al sistema.}
    }
}
\subsection{Spazi di Hilbert}
Si descrive ora più precisamente cosa si intende per spazi di Hilbert, definendo anche
più rigorosamente gli stati normalizzabili "alla Dirac".

Si consideri uno spazio vettoriale $ E $ definito sul campo complesso $ \mathbb{C} $:
esso si dice unitario se è dotato di un prodotto scalare $(\cdot,\cdot)$; ogni spazio
unitario è anche metrico rispetto alla distanza $ d(a,b) = \sqrt{(a-b,a-b)} $ e questa
distanza induce una topologia simile a quella di $ \mathbb{R}^n $.

Con questa topologia si può definire l'usuale nozione di limite di successioni e si 
può dimostrare il seguente teorema

\noindent\fbox{%
    \parbox{\textwidth}{%

	\textbf{Teorema}, continuità del prodotto scalare:
	$$ \lim\limits_{n\to\infty} a_n = l \rightarrow \forall b\in E, \lim\limits_{n\to\infty} (b,a_n) = (b,l)$$
	\textbf{Dimostrazione:} Preso $b\in E$, si consideri $(b,a_n) - (b,l) = (b, a_n - l) $:	per la disuguaglianza di Schwartz, $ || (b,a_n-l) ||^2 \leq(b,b)(a_n-l,a_n-l)$,
	ma $\lim\limits_{n\to\infty} (a_n-l,a_n-l) = 0$,
    
	dunque $\lim\limits_{n\to\infty} (b,a_n-l)=0$, \textbf{QED}. 

    }%
}


Un altro concetto molto utile per la formulazione della Meccanica quantistica è quello di
spazio duale:

\textbf{Def.} Dato $ E $ spazio vettoriale, il suo \textit{duale }è
$ E^* := \left\{ \hat{b}: E\to\mathbb{C}\: /\: \hat{b} \text{ è lineare e continuo} \right\} $

Poiché $ E $ è uno spazio unitario, gli elementi di $ E^* $ si possono identificare con
$ \hat{b} = (b,\cdot) $,

dove $ \hat{b}(a) = (b,a)\: \forall a\in E $.

Da qui in poi si userà la \textbf{notazione di Dirac}, indicando gli elementi di
$ E^* $ con $ \hat{b} := \bra{b} $ ("bra"), quelli di $ E $ con $ a = \ket{a} $ ("ket") e
i prodotti scalari con $(a,b):=\bra{a}\ket{b}$ ("bra-(c)-ket").

Quanto descritto finora vale per spazi vettoriali di dimensione sia finita, sia infinita,
ma per procedere a trattarli ugualmente bisogna definire una classe \textit{well-behaved}
di spazi a dimensione infinita.

\textbf{Def.} Uno spazio vettoriale si dice \textit{separabile} se ammette sottoinsiemi 
$ \left\{ \ket{e_i}\right\}$ completi ed enumerabili\footnote{Non è la più generale
definizione di separabilità, ma in questo contesto basta.}, che può essere scelta
ortonormale.
Questo garantisce che

$$ \forall \ket{v}\in E,\: \exists c_i \in \mathbb{C} \:\text{t.c.}\:
\sum_{i}^{\infty} |c_i| = \bra{v}\ket{v} \leftrightarrow
\sum_{i}^{\infty}c_i\ket{e_i} = \ket{v} \leftrightarrow
\sum_{i}^{\infty}\ket{e_i}\bra{e_i} = \mathbb{I}
$$
che sono relazioni banalmente verificate per spazi a dimensione finita.

\textbf{Def.} Una successione $ \ket{a}_n $ si dice di Cauchy se
$$ \forall \epsilon>0\;\exists n_0\in\mathbb{N} \text{ t.c } \forall n,m>n_0,\:
d(\ket{a}_n - \ket{a}_m)<\epsilon$$
\textbf{Oss.} Ogni successione convergente è una successione di Cauchy, ma il viceversa
non è sempre garantito.

\textbf{Def.} Uno spazio metrico si dice \textit{completo} se ogni successione di Cauchy
è una successione convergente.

Si può finalmente dare la definizione di spazio di Hilbert

\textbf{Def.} Uno spazio unitario e completo si dice \textbf{spazio di Hilbert}.

L'esempio principale di spazio di Hilbert è $ L^2(a,b) $, cioè l'insieme delle funzioni
quadrato-sommabili sull'intervallo $ (a,b)\subseteq\mathbb{R} $, dotato del prodotto
scalare $ \bra{f}\ket{g} :=\int_{a}^{b}dxf^*(x)g(x) $.

Già durante il corso di meccanica quantistica I si sono incontrate però delle soluzioni
dell'equazione di Schrödinger \textit{non} appartenenti a $ \mathcal{H} $ utilizzate
però come "base generalizzata" dello spazio: si cerca ora di formalizzare questo
processo.

\subsection{Distribuzioni temperate e spazi di Hilbert}

\textbf{Def.} Lo spazio delle funzioni di prova è
$$ S := \left\{ f(x) / f\in C^{\infty} \land \lim\limits_{|x|\to\infty}
x^p \dv[n]{f}{x}=0,\; \forall p,n\in\mathbb{N}\right\} $$

Si dice che $ f(x) $ è rapidamente decrescente, cioè che essa e tutte le sue derivate
decrescano a infinito più rapidamente di ogni potenza.

\textbf{Oss.} Si può dimostrare che $ S $ è denso in $L^2(\mathbb{R})$, cioè che
$ \forall$ intorno di $\ket{f}\in L^2(\mathbb{R})$, esiste almeno una $\ket{g}\in S$. 

\textbf{Def.} Lo spazio delle distribuzioni temperate è
$$ S':=\left\{ T:S\to\mathbb{C}/T\text{ è lineare e continua} \right\} $$

Il numero $ \bra{T}\ket{g}\in\mathbb{C} $ denota l'azione di $ \bra{T}\in S' $ su $ \ket{g}\in S $.

L'azione di $ T $ si comporta come un prodotto scalare, infatti
\begin{itemize}
    \item È lineare per definizione: $\bra{T}\ket{f+g}=\bra{T}\ket{f} + \bra{T}\ket(g)$;
    \item È continua per definizione: 
	$ \bra{T}\ket{\lim\limits_{n\to\infty}g_n} =\lim\limits_{n\to\infty}
	\bra{T}\ket{g_n} $, con il limite di funzioni di $ S $ a primo membro. 
\end{itemize}


L'insieme $ S' $ contiene anche funzioni nel senso usuale: ogni funzione $ T(x) $ 
localmente integrabile e a crescenza algebrica finita (cioè che a infinito al più
diverge come un polinomio) definisce una distribuzione $ \bra{T} $ (detta
\textit{regolare}) definita da $\bra{T}\ket{g} := \int_{-\infty}^{\infty}dxT^*(x)g(x) $\footnote
{Questo processo ricorda l'identificazione degli elementi dello spazio duale}.

L'esempio del contrario è la famosa $ \delta $ di Dirac: è infatti una distribuzione
definita da $ \bra{\delta}\ket{f}:=f(0)$, $ \forall f\in S $; l'espressione
$ \int_{-\infty}^{\infty}dx\delta(x) f(x) = f(0) $ è un abuso di notazione, perché non 
esiste alcuna funzione $ \delta(x) $.

Come si è visto nel corso di Metodi I, si può però aggirare il problema delle
distribuzioni non regolari notando che vale il seguente teorema:

\noindent\fbox{%
    \parbox{\textwidth}{%
	\textbf{Teorema:} $ \forall T\in S' $, $ \exists \left\{ T_n \right\}$,
	successione di distribuzioni regolari, tale che
	$ \lim\limits_{n\to\infty}^{deb.} T_n = T $  
    }%
}

Ricordando la definizione di limite debole:

\textbf{Def.} Data $ T\in S' $ e $\left\{ T_n \right\}$ successione di distribuzioni
regolari, si dice che 
$$ \lim\limits_{n\to\infty}^{deb.} T_n = T \leftrightarrow 
\forall g\in S, \lim\limits_{n\to\infty} \bra{T_n}\ket{g} = \bra{T}\ket{g}$$

In questo modo, si possono utilizzare successioni di distribuzioni regolari per
trattare di distribuzioni più generali (si veda la trattazione della $ \delta $ di 
Dirac nel corso di Metodi I).

In conclusione, per trattare gli stati non normalizzabili è necessario introdurre sia
$ S $, sia $ S' $, dando la seguente definizione

\textbf{Def.} La terna (di Gel'Fand) $ S\subset\mathcal{H}\subset S' $, con $S$ denso in
$ \mathcal{H} $ e $ S' $ spazio delle distribuzioni temperate su $ S$, si dice 
\textbf{spazio di Hilbert equipaggiato}.
\newpage

\section{Postulato 0.1 della Meccanica Quantistica}

\noindent\fbox{%
    \parbox{\textwidth}{%
    \textbf{Ogni quantità misurabile del sistema fisico è associata a un'osservabile,
    cioè un operatore autoaggiunto con dominio denso in $\mathcal{H}$.
    Gli autovettori (propri o generalizzati) dell'osservabile formano una base
    ortonormale per $ \mathcal{H} $ }
    }%
}

Per comprendere questa definizione e formalizzarla è necessario definire alcuni concetti sugli
operatori e sui loro autovettori.

\subsection{Osservabili}

Un operatore quantistico è una funzione $ A: \mathcal{D}_A\subset\mathcal{H}\to\mathcal{H} $
continua, dove $ \mathcal{D}_A $ è il dominio dell'operatore $ A $.

\textbf{Def.} L'operatore aggiunto di $ A $ è l'unico operatore $ A^{\dagger} $ tale che
$ \forall \ket{v}\in\mathcal{D}_A,\; \forall \ket{u}\in\mathcal{D}_{A^{\dagger}}$ vale

$$ \bra{u} \left( A\ket{v} \right) = \left[ \bra{v} \left( A^{\dagger}\ket{u} \right) \right]^* $$ 
 
\textbf{Def.} Un operatore $ A $ si dice autoaggiunto se $ A = A^{\dagger} $ e
$\mathcal{D}_A = \mathcal{D}_{A^{\dagger}}$

\textbf{Oss.} È una definizione più stringente della hermiticità, come è chiaro dal seguente
esempio:

\textbf{Esempio} Si consideri $ A = -i\dv{x} $ nello spazio $ \mathcal{H} = L^2(a,b) $: 
questo operatore è sempre hermitiano, ma non è sempre autoaggiunto, infatti
$$ \bra{f}A\ket{g} = \int_{a}^{b}dxf^*(x)\left(-i\dv{x}\right)g(x) =
-i\eval{g(x)f(x)}_a^b + \int_{a}^{b}dx\left(i\dv{x}\right)f^*(x)g(x)$$
è pari a $ \left[ \bra{g}A^{\dagger}\ket{f} \right]^* $ solo se sia $ f(x) $ sia $ g(x) $
sono periodiche su $ (a,b) $.

L'operatore è dunque autoaggiunto\footnote{Questo esempio non è precisissimo, per una trattazione
perfetta si vedano le dispense del prof. Sciuto, pp. 156-157.} in
$ \mathcal{D}_A = \left\{ f:(a,b)\to\mathbb{C} \text{ periodiche} \right\}$, per $ (a,b)$ finito.

Il concetto di osservabile si può estendere anche alle distribuzioni temperate
\footnote{Anche se esse non rappresentano propriamente stati fisici, come si vedrà
in seguito}:

\textbf{Def.} Dato $ A$ con $ \mathcal{D}_A\subset S $, si definisce l'azione di $ A $ su 
$ \bra{T} \in S'$ come $ A: S'\to S' $ tale che
$$ \left(\bra{T}A\right)\ket{g} := \bra{T}\left(A\ket{g}\right),\; \forall \ket{g}\in S \text{ t.c } A\ket{g}\in S$$
\textbf{Oss.} Se si considera il caso in cui $ T$ è una funzione, la definizione è coerente con
quella di prodotto scalare tra \textit{bra} e \textit{ket}.

\textbf{Oss.} Da questa definizione si può anche definire l'espressione $\bra{g}A\ket{T}$
nel modo seguente: $ \left[\bra{T}A^{\dagger}\ket{g}\right]^* := \bra{g}A\ket{T} $ 
